% BHU PYQ -- History: History of India (1740 - 1857) (2025-26 NEP)
\documentclass[11pt,a4paper]{article}
\usepackage[a4paper,top=2.5cm,bottom=2.5cm,left=2.8cm,right=2.8cm,headheight=14pt]{geometry}
\usepackage{amsmath,amssymb}
\usepackage{fontspec}
\setmainfont{Kohinoor Devanagari}
\usepackage{microtype,fancyhdr,enumitem,hyperref,lastpage,parskip}

\hypersetup{colorlinks=true,urlcolor=black,linkcolor=black,pdftitle={HISMJ31: History of India (1740 - 1857) -- BHU NEP 2025-26}}

\pagestyle{fancy}\fancyhf{}
\renewcommand{\headrulewidth}{0.4pt}\renewcommand{\footrulewidth}{0.4pt}
\fancyhead[L]{\small \textit{Banaras Hindu University}}
\fancyhead[R]{\small \textit{HISMJ31: History of India (1740 - 1857) (2025-26)}}
\fancyfoot[C]{\small Page \thepage\ of \pageref{LastPage}}

\newlist{parts}{enumerate}{2}
\setlist[parts,1]{label=(\alph*),leftmargin=*,itemsep=4pt,topsep=3pt}
\setlist[parts,2]{label=(\roman*),leftmargin=*,itemsep=2pt,topsep=2pt}

\newcommand{\pts}[1]{\hfill \textit{[#1]}}

\begin{document}

\begin{center}
  {\large\bfseries BANARAS HINDU UNIVERSITY}\\[2pt]
  {\normalsize B. A. (Hons.) Semester III Examination, 2025-26 (NEP)}\\[6pt]
  \rule{0.8\linewidth}{0.5pt}\\[6pt]
  {\large\bfseries HISTORY}\\[4pt]
  {\large\bfseries Paper Code: HISMJ-31 : History of India (1740 -- 1857)}\\[4pt]
  {\normalsize Time: 2:30 Hours \hfill Full Marks: 70}\\[2pt]
  {\small REF NO. 2025-26/D-III/141}
\end{center}
\rule{\linewidth}{0.4pt}
\smallskip

\noindent \textit{Instructions: This question paper comprises three Sections A, B and C. All questions are compulsory. Write your Roll No.\ at the top immediately on receipt of this question paper.}

\smallskip
\noindent \textit{इस प्रश्न-पत्र में तीन खण्ड अ, ब तथा स हैं। सभी प्रश्न अनिवार्य हैं।}

\bigskip

\section*{SECTION -- A / खण्ड -- अ}
\noindent \textit{Note: Objective / Very Short Answer Questions. Answer each in 15--20 words.}\pts{5 $\times$ 2 = 10}

\noindent \textit{वस्तुनिष्ठ / अति लघु उत्तरीय प्रश्न। प्रत्येक प्रश्न का उत्तर 15--20 शब्दों में दीजिए।}

\begin{enumerate}[label=\textbf{\arabic*.},leftmargin=*,itemsep=8pt]
  \item 
  \begin{parts}
    \item Write about Robert Clive.\pts{2}
    
    क्लाइव के विषय में लिखिए।

    \item Who was Mir Qasim?\pts{2}
    
    मीर कासिम कौन था?

    \item Write two provisions of Pitt's India Act of 1784.\pts{2}
    
    पिट्स इण्डिया एक्ट (1784) के दो प्रावधान लिखिए।

    \item What do you mean by the Dual Government system in Bengal?\pts{2}
    
    द्वैध शासन प्रणाली से आप क्या समझते हैं?

    \item What was `Dastak'?\pts{2}
    
    `दस्तक' क्या था? (Free trade permit issued to British East India Company)
  \end{parts}
\end{enumerate}

\bigskip

\section*{SECTION -- B / खण्ड -- ब}
\noindent \textit{Note: Short Answer Type Questions. Answer each in about 250 words.}\pts{3 $\times$ 10 = 30}

\noindent \textit{लघु उत्तरीय प्रश्न। प्रत्येक प्रश्न का उत्तर लगभग 250 शब्दों में दीजिए।}

\begin{enumerate}[label=\textbf{\arabic*.},leftmargin=*,start=2,itemsep=12pt]

  \item Discuss the causes and consequences of the Battle of Plassey (1757).\pts{10}
  
  प्लासी के युद्ध (1757) के कारणों एवं परिणामों की चर्चा कीजिए।

  \begin{center}\textbf{OR / अथवा}\end{center}

  Discuss in detail the contribution of Raja Rammohan Roy and Brahmo Samaj in the social reform movement of 19th century India.\pts{10}
  
  19वीं शताब्दी के भारतीय सामाजिक सुधार आन्दोलन में राजा राममोहन राय एवं ब्रह्म समाज के योगदान की विस्तार से चर्चा कीजिए।

  \item Mention the circumstances that led to the Battle of Buxar (1764) and explain its significance.\pts{10}
  
  बक्सर के युद्ध (1764) के लिए उत्तरदायी परिस्थितियों का उल्लेख कीजिए तथा इसके महत्त्व की व्याख्या कीजिए।

  \begin{center}\textbf{OR / अथवा}\end{center}

  Outline the main features of the Permanent Settlement of Bengal (1793).\pts{10}
  
  बंगाल के स्थायी बन्दोबस्त (1793) की प्रमुख विशेषताओं को रेखांकित कीजिए।

  \item Write a brief essay on the background and provisions of the Regulating Act of 1773.\pts{10}
  
  1773 के रेग्युलेटिंग एक्ट की पृष्ठभूमि तथा प्रावधानों पर एक संक्षिप्त निबन्ध लिखिए।

  \begin{center}\textbf{OR / अथवा}\end{center}

  Discuss the Subsidiary Alliance system propounded by Lord Wellesley.\pts{10}
  
  लॉर्ड वेलेज़ली द्वारा प्रतिपादित सहायक संधि प्रणाली की विवेचना कीजिए।

\end{enumerate}

\bigskip

\section*{SECTION -- C / खण्ड -- स}
\noindent \textit{Note: Long Answer Type Questions. Answer each in about 500 words.}\pts{2 $\times$ 15 = 30}

\noindent \textit{दीर्घ उत्तरीय प्रश्न। प्रत्येक प्रश्न का उत्तर लगभग 500 शब्दों में दीजिए।}

\begin{enumerate}[label=\textbf{\arabic*.},leftmargin=*,start=5,itemsep=14pt]

  \item Write an essay on the nature, background and causes of the Revolt of 1857.\pts{15}
  
  1857 की क्रान्ति (विद्रोह) की प्रकृति, पृष्ठभूमि एवं कारणों पर एक निबन्ध लिखिए।

  \begin{center}\textbf{OR / अथवा}\end{center}

  Trace the various phases of British educational policy in India from 1835 to 1857.\pts{15}
  
  1835 से 1857 तक ब्रिटिश शिक्षा नीति के विभिन्न चरणों का रेखांकन कीजिए।

  \item Critically explain the social policy of the British East India Company between 1829 and 1857.\pts{15}
  
  1829 से 1857 के मध्य ब्रिटिश ईस्ट इण्डिया कम्पनी की सामाजिक नीति की समालोचनात्मक व्याख्या कीजिए।

  \begin{center}\textbf{OR / अथवा}\end{center}

  Which economic policies of the British East India Company between 1813 and 1857 led to famines in India? Explain.\pts{15}
  
  1813 से 1857 के मध्य ब्रिटिश ईस्ट इण्डिया कम्पनी की किन आर्थिक नीतियों के कारण भारत में अकाल पड़े? व्याख्या कीजिए।

\end{enumerate}

\end{document}
